\chapter{46}

\section{Dienstag, 21. August}



Der ordentliche Arbeitsmorgen war bereits ein paar Stunden alt.
Mantovani trat gerade durch die Tür auf den Flur der Quästur. Frisch
geduscht und kaffeegestärkt machte er den Eindruck, die schlafarme Nacht
weggesteckt zu haben. Drei seiner Ermittler fand er stehend vor dem
Kaffeeautomaten. Der eine hob den Daumen.

«Wir haben was, Commissario.»

Mantovani bedeutete ihm und Melik Sönmez, ihm in sein Büro zu folgen.
Die beiden stürzten ihre Espressi hinunter. Sie warfen ihre Becher
während des Gehens in einen Abfallkorb und folgten ihm.

Zu dritt sassen sie um Mantovanis Bürotisch. Dieser studierte seine
Handy-Notizen. Mit sachlicher Genauigkeit fasste er den beiden die
Diagnose des Arztes zu Giorgias Zustand zusammen.

Ermittler Federico Rossi schüttelte den Kopf.

Mantovani schaute auf.

«Eine quere Geschichte,» begann Sönmez, er runzelte die Stirn, «also,
wir haben oben auf dem Platz, direkt an der Felskante und in
unmittelbarer Nähe, einiges gefunden. Da muss etwas stattgefunden haben,
eine Party vielleicht…der Platz war leer, aber die Glut der
Feuerstelle etwa 50 Meter vom Felsabgang entfernt, war noch warm. Einen
halben Meter entfernt haben wir eine kleine Damen-Handtasche gefunden.
Sie gehört laut Ausweis in einem darin liegenden Geldbeutel, einer
\emph{Giorgia Sala}, wohnhaft in Mailand, \emph{Via Rosaria 39}, und,
wir haben Fingerabdrücke am Taschenhenkel gefunden. Die konnten bereits
identifiziert werden. Die stammen von einem,» Sönmez senkte den Blick
und las von seinen Notizen ab, «\emph{Bastiano Gabriele di Martino},»
auch ein Milanese, ist im System, vorbestraft wegen verschiedener
Vergehen, Drogenhandel im kleinen Stil, aber immerhin. Sichergestellt
haben wir auch frische Fussspuren, direkt an der Kante und auffallend
gut sichtbar in den trocken-sandigen Boden vertieft. Die sind bereits
bei der Spurensicherung, und,» Sönmez atmete kurz durch, «zwei Kondome,
rund 10 Meter von der Tasche entfernt unter einem Gebüsch, plus eins
etwas weiter entfernt auf dem schmalen Pfad Richtung Osten, den Hügel
runter. Alle drei gleiche Marke, die zwei im Gebüsch benutzt, das dritte
auf dem Weg unbenutzt verschlossen, die sind bereits im Labor, um…»

Bevor er weitersprechen konnte, unterbrach ihn Mantovani.

«Sala? Wie sagtest du, ist die Adresse dieser \emph{Giorgia Sala}?»

«Via Rosaria 39.»

Sönmez legte seine Papiere auf den Tisch.

«Das ist doch die Adresse der \emph{Tanzschule Sala}?»

Mantovanis Frage blieb zwischen ihnen liegen. Rossi und Sönmez sahen
sich schulterzuckend an.

«Ich habe mit Anna vergangenen Winter einen Tanzkurs gemacht bei den
Salas. Giorgia Sala könnte entsprechend ihres Jahrgangs ihre Tochter
sein, prüft das, bitte.»

Mantovani sah betroffen aus.

«Und,» fuhr er fort, «\emph{di Martino} heisst dieser Mann, den ihr in
der Datenbank gefunden habt?»

Sömnez bejahte.

«Klärt erstmal ab, ob der Mann…, wie alt ist der?»

«Jahrgang, Moment,» Sömnez warf einen Blick auf seine Notizen, «Jahrgang
91.»

Er blickte interessiert zu Mantovani hoch.

Mantovanis Augen wanderten zur Decke, wie immer, wenn er gedanklich
entschwand.

«Hm, dann klärt bitte ab, ob \emph{Bastiano di Martino} der Sohn der
einflussreichen Mailänder Familie \emph{Bernardo di Martino} ist.
Ehrlich gesagt, hoff ich das nicht.»

Seine Stimme hörte sich nun brummend an.

«Wenn, dann kann das ja heiter werden.»

Auch der Name \emph{di Martino} war den beiden unbekannt.

«Und bitte, spürt Bastiano di Martino auf, Wohnort, Arbeitsort, und
sprecht eine Vorladung aus für heute, späten Nachmittag, siebzehn Uhr.
Die Schuhspuren sind bis dahin noch nicht bestimmt, aber die
Fingerabdrücke auf der Tasche reichen aus. Lasst auch die Kondome auf
Fingerabdrücke untersuchen. Vielleicht schaffen sie’s bis um fünf. Eine
Übereinstimmung mit diesem \emph{di Martino} könnte interessant sein.
Ich werde ihn persönlich vernehmen.»

«Geht klar, Commissario.»

Die beiden Ermittler wollten gerade aufstehen, als Mantovani sie bat, zu
warten.

«Über den Verrückten, der Cristina mit seinem Auftritt überfallen hat,
wisst ihr auch schon was?»

«Leider nein,» Rossi schüttelte genervt den Kopf, «für die kurze Zeit
des Auftritts des Mannes ist Cristinas Beschreibung ziemlich gut, aber
wir sind noch nicht dazugekommen, weil…»

«Versteh ich, klar, bitte beschäftigt euch damit. Es liegt nah, dass der
mit dem Absturz der Frau etwas zu tun hat. Findet raus, wer er ist und
warum er hier bei uns aufgetaucht ist. Dazu braucht ihr Verstärkung.
Emma Riva wird euch unterstützen. Setzt euch als Erstes mit ihr
zusammen, okay?»

«Ja, Commissario.»

Rossi sammelte sich.

«Vorladung di Martino, Auswertung der weiteren Spuren und Absprache mit
Emma.»

Die beiden Polizisten verliessen den Raum.

Mantovani stützte sein Gesicht auf den Händen ab, zog die Unterlippe ein
und atmete tief aus. Wenn es sich tatsächlich um den Sohn \emph{der} di
Martino handelte, kamen beschwerliche Ermittlungen auf sie zu. Der
Einfluss dieser betuchten Familie war in Mailand beängstigend gross.
Mantovani wusste, dass sie in etliche unsaubere Geschäfte verwickelt
waren. Und auch, dass sie mit Hilfe ihrer Vernetzung mit andern
Mailänder Grössen, bis ganz oben, kaum je belangt werden konnten. An
denen hatten sich schon andere die Zähne ausgebissen.
